\section{Final Test (2022--2025)}
\label{sec:final}

\begin{table*}[t]
\caption{Final test, 2022--2025 (997 signal dates), run once with the configuration frozen beforehand. The equal-weight tradable universe earned 11.2\%/yr and the S\&P~500 11.1\%/yr; returns are net of 5\,bp per side. P(IR$>$0) is the probabilistic Sharpe ratio of the active return against zero; the last column repeats the development IR.}
\label{tab:final}
\footnotesize
\begin{tabular}{@{}lrrrrrrr@{}}
\toprule
Model & \ric{} ($t$) & Return (\%/yr) & Net active (\%/yr) & IR (active $t$) & IR at 10\,bp & P(IR$>$0) & IR, development \\
\midrule
\tablefile{tables/final_results.tex}
\bottomrule
\end{tabular}
\end{table*}

\begin{table*}[t]
\caption{The eight paired comparisons declared before the final test, 2022--2025. No interval excludes zero.}
\label{tab:final-gaps}
\footnotesize
\begin{tabular}{@{}lll@{}}
\toprule
Comparison (A minus B) & \ric{} difference [95\% CI] & Net return difference, \%/yr [95\% CI] \\
\midrule
\tablefile{tables/final_gaps.tex}
\bottomrule
\end{tabular}
\end{table*}

Before opening the final test, we recorded in a configuration file the primary model, six comparators, eight paired comparisons to report, a hash of the code and the number of development trials.
Each model followed its development procedure, with LightGBM retrained every year and TimesFM before 2022 and 2025, and the results in \cref{tab:final} are reported as they came out.

\begin{table}[t]
\caption{Net active return by year on the final test, \%/yr at 5\,bp.}
\label{tab:final-years}
\footnotesize
\setlength{\tabcolsep}{4pt}
\begin{tabular}{@{}lrrrr@{}}
\toprule
Model & 2022 & 2023 & 2024 & 2025 \\
\midrule
\tablefile{tables/final_by_year.tex}
\bottomrule
\end{tabular}
\end{table}

\paragraph{The primary model replicates its development performance.}
It earned 17.8\%/yr, +6.7\%/yr above the equal-weight universe after costs, with an IR of 0.60, against +6.8\%/yr and 0.64 in development.
It keeps +4.3\%/yr at 10\,bp, its \ric{} is 0.0263 ($t=2.0$), and its largest drawdown is $-25\%$.
Four years give little power, though: an IR of 0.6 implies an active $t$ of only 1.2, and the active return swings from year to year (+17.1\%, $-6.9$\%, +4.8\% and +11.8\% from 2022 to 2025; \cref{tab:final-years}).
None of the eight declared comparisons excludes zero (\cref{tab:final-gaps}); for instance, the primary model beats reversal by +3.2\%/yr with an interval of [$-10.3$, +14.1].

\paragraph{Four-year rankings are noisy.}
Momentum has the best final-test IR (0.69), driven by an active return of +32.4\% in 2024, after a development IR of only 0.09, while reversal falls to 0.24.
Switching to momentum now would repeat the selection error that the frozen protocol is meant to prevent.

\paragraph{TimesFM on unseen data.}
Zero-shot TimesFM again shows no skill (\ric{} 0.0053, IR 0.05).
Fine-tuned TimesFM does better than in development (IR 0.49), and its \ric{} of 0.0194 is the only TimesFM result whose bootstrap interval excludes zero on unseen data ([0.0017, 0.0397]); its gain over zero-shot (+4.9\%/yr) is not significant, and it stays below the primary model.
The combination with reversal falls from an IR of 0.80 to 0.42, consistent with its dependence on the 2008--2010 crisis.
